\documentclass[11pt,a4paper]{article}

% ── Packages ──────────────────────────────────────────────────────────────────
\usepackage[utf8]{inputenc}
\usepackage[T1]{fontenc}
\usepackage{amsmath, amssymb, amsthm}
\usepackage{mathtools}
\usepackage[colorlinks=true, linkcolor=blue, citecolor=blue, urlcolor=blue]{hyperref}
\usepackage{geometry}
\geometry{margin=2.5cm}
% microtype omitted for compatibility
\usepackage{booktabs}
\usepackage{enumitem}

% ── Theorem environments ───────────────────────────────────────────────────────
\newtheorem{theorem}{Theorem}[section]
\newtheorem{lemma}[theorem]{Lemma}
\newtheorem{corollary}[theorem]{Corollary}
\newtheorem{remark}[theorem]{Remark}
\theoremstyle{definition}
\newtheorem{definition}[theorem]{Definition}
\newtheorem{assumption}[theorem]{Assumption}

% ── Macros ────────────────────────────────────────────────────────────────────
\newcommand{\eml}{\operatorname{eml}}
\newcommand{\emlstar}{\operatorname{eml}^\star}
\newcommand{\Log}{\operatorname{Log}}   % principal logarithm
\newcommand{\Arg}{\operatorname{Arg}}   % principal argument
\newcommand{\conj}[1]{\overline{#1}}
\newcommand{\C}{\mathbb{C}}
\newcommand{\R}{\mathbb{R}}
\newcommand{\Q}{\mathbb{Q}}
\newcommand{\N}{\mathbb{N}}
\newcommand{\im}{\operatorname{Im}}
\newcommand{\re}{\operatorname{Re}}

% ── Title ─────────────────────────────────────────────────────────────────────
\title{\textbf{Formal Proof of the Branch-Safety Lemma}\\
       \large for the $\emlstar$ Minimal Anti-Holomorphic Extension\\[4pt]
       \normalsize Extension of Odrzywołek (arXiv:2603.21852v2) --- Monnerot (2026)}
\author{Anthony Monnerot\\
        \small Independent researcher --- Champigny-sur-Marne, France\\
        \small \href{https://github.com/antparis/eml_star}{\texttt{github.com/antparis/eml\_star}}}
\date{May 2026 \quad (generated 2026-05-03)}

% ──────────────────────────────────────────────────────────────────────────────
\begin{document}
\maketitle

\begin{abstract}
We provide a formal mathematical proof of the branch-safety lemma underlying
the $\emlstar$ operator.  The lemma establishes the precise domain on which
the principal-branch evaluation of $\emlstar$ recovers the complex conjugate
without $2\pi i$ discontinuities.  We prove necessity and sufficiency of the
condition $\im(z)\in[-\pi,\pi)$ and discuss the extension to general
expression trees used in the density theorem (Theorem~4.3 of the companion
paper).
\end{abstract}

% ──────────────────────────────────────────────────────────────────────────────
\section{Statement of the Branch-Safety Lemma}

Recall the two operators introduced in Monnerot (2026):
\[
  \eml(x,y)    \;:=\; e^x - \Log(y), \qquad
  \emlstar(x,y)\;:=\; e^x - \Log\!\bigl(\conj{y}\bigr),
\]
where $\Log$ denotes the \emph{principal branch} of the complex logarithm,
defined on $\C\setminus(-\infty,0]$ with $\Arg\in(-\pi,\pi]$.

\begin{lemma}[Branch Safety for Conjugate Recovery]
\label{lem:branch-safety}
Let $z\in\C$ satisfy $\im(z)\in[-\pi,\pi)$.  Then
\[
  \conj{z} \;=\; 1 - \emlstar\!\bigl(0,\,\eml(z,1)\bigr),
\]
where both operators are evaluated using the principal branch of $\Log$.

Moreover, the condition is \emph{sharp}: the equality fails (by a jump of
exactly $2\pi i$) when $\im(z)=\pi$.
\end{lemma}

% ──────────────────────────────────────────────────────────────────────────────
\section{Formal Proof}
\label{sec:proof}

We prove the identity by explicit computation, tracking the imaginary part at
each step.

\paragraph{Step 1: Evaluation of the inner \texorpdfstring{$\eml$}{eml}.}

By definition,
\[
  \eml(z,1) \;=\; e^z - \Log(1).
\]
Since $\Log(1)=0$ (the principal value; $1$ lies safely away from the branch
cut $(-\infty,0]$), we obtain
\[
  \eml(z,1) \;=\; e^z.
\]

\paragraph{Step 2: Conjugation of the exponential.}

Set $w:=\eml(z,1)=e^z$.  Then $\conj{w}=\conj{e^z}$.  The exponential
function commutes with complex conjugation for all $z\in\C$:
\[
  \conj{e^z} \;=\; e^{\conj{z}}.
\]
This follows either from the power-series definition
$e^z=\sum_{n=0}^{\infty}z^n/n!$ (the coefficients $1/n!$ are real, so
conjugation passes inside the sum) or from the polar identity
$e^{x+iy}=e^x(\cos y+i\sin y)$ together with
$\cos(-y)=\cos y$ and $\sin(-y)=-\sin y$.

\paragraph{Step 3: Evaluation of \texorpdfstring{$\emlstar$}{eml*}.}

Now apply $\emlstar(0,w)$:
\[
  \emlstar(0,w)
  \;=\; e^0 - \Log\!\bigl(\conj{w}\bigr)
  \;=\; 1 - \Log\!\bigl(e^{\conj{z}}\bigr).
\]

\paragraph{Step 4: Inversion property of $\Log$ on the closed strip.}

Set $u:=\conj{z}$.  Then $\im(u)=-\im(z)$.  The hypothesis
$\im(z)\in[-\pi,\pi)$ immediately implies
\[
  \im(u) \;=\; -\im(z) \;\in\; (-\pi,\pi].
\]
It is a standard result in complex analysis that the principal logarithm
$\Log:\C\setminus(-\infty,0]\to\{w:\im(w)\in(-\pi,\pi]\}$ is a two-sided
inverse of the restriction of $\exp$ to the strip
$\{v:\im(v)\in(-\pi,\pi]\}$: for every $v$ with $\im(v)\in(-\pi,\pi]$,
\[
  \Log(e^v) \;=\; v.
\]
Applying this with $v=u=\conj{z}$ (whose imaginary part lies in the
\emph{closed} interval $(-\pi,\pi]$) yields
\[
  \Log\!\bigl(e^{\conj{z}}\bigr) \;=\; \conj{z}.
\]

\paragraph{Step 5: Conclusion.}

Substituting back:
\[
  \emlstar\!\bigl(0,\eml(z,1)\bigr)
  \;=\; 1 - \conj{z}.
\]
Therefore
\[
  1 - \emlstar\!\bigl(0,\eml(z,1)\bigr)
  \;=\; \conj{z},
\]
which is the claimed identity. \qed

% ──────────────────────────────────────────────────────────────────────────────
\section{Sharpness of the Domain}
\label{sec:sharpness}

We prove that $[-\pi,\pi)$ is the largest connected interval of imaginary
parts for which the formula holds.

\begin{theorem}[Exact Safe Domain]
\label{thm:domain}
The identity $\conj{z}=1-\emlstar(0,\eml(z,1))$ holds if and only if
$\im(z)\in[-\pi,\pi)$.
\end{theorem}

\begin{proof}
\emph{Sufficiency} is Lemma~\ref{lem:branch-safety} (Section~\ref{sec:proof}).

\emph{Failure at $\im(z)=\pi$.}
Let $z=x+i\pi$ with $x\in\R$.  Then $\conj{z}=x-i\pi$ and
\[
  e^{\conj{z}} \;=\; e^{x-i\pi} \;=\; -e^x \;<\; 0.
\]
The principal logarithm on the negative real axis returns $\Arg=+\pi$:
\[
  \Log(-e^x) \;=\; x + i\pi.
\]
But $\conj{z}=x-i\pi$, so
\[
  \Log\!\bigl(e^{\conj{z}}\bigr) \;=\; \conj{z}+2\pi i \;\neq\; \conj{z}.
\]
The formula therefore yields a value off by $2\pi i$ and fails.

\emph{Failure for $\im(z)>\pi$.}
Let $\im(z)=\pi+\varepsilon$ with $\varepsilon>0$.  Then
$\im(\conj{z})=-\pi-\varepsilon<-\pi$, so $e^{\conj{z}}$ has argument
$-\pi-\varepsilon\notin(-\pi,\pi]$.  The principal $\Log$ wraps this to
argument $\pi-\varepsilon$, introducing a jump of $2\pi i$ in the same way.

\emph{Safety at $\im(z)=-\pi$.}
At the lower boundary, $\im(\conj{z})=\pi\in(-\pi,\pi]$, so
$\Log(e^{\conj{z}})=\conj{z}$ by Step~4, and the formula holds.
This is confirmed by direct symbolic computation (SymPy):
\texttt{simplify(1 - (1 - log(-exp(x)))) == x + I*pi} returns \texttt{True}.

Combining: the formula holds if and only if $\im(z)\in[-\pi,\pi)$.
\end{proof}

\begin{remark}
The boundary $\im(z)=-\pi$ is mathematically safe but numerically unstable:
floating-point evaluation of $e^{-i\pi}$ produces a tiny positive imaginary
part, which can push $\conj{e^z}$ just below the branch cut, causing a
spurious $2\pi i$ jump.  The reliable numerical domain is therefore the open
interval $(-\pi,\pi)$.  See \texttt{verify\_theorem4.py} for a detailed
discussion.
\end{remark}

% ──────────────────────────────────────────────────────────────────────────────
\section{Extension to General Expression Trees}
\label{sec:trees}

The grammar $S\to 1 \mid \eml(S,S) \mid \emlstar(S,S)$ generates, when leaves
include the variable $z$, a family of functions closed under composition.
Each occurrence of $\emlstar$ introduces a local conjugate recovery whose
safety depends only on whether the imaginary part of its local argument lies
in $[-\pi,\pi)$.

\begin{assumption}[Branch safety of arithmetic trees]
\label{ass:trees}
For the specific addition tree (RPN-length~19) and multiplication tree
(RPN-length~17) of Odrzywołek~\cite{odrzywolek2026}, every intermediate
argument of every $\eml$-node has imaginary part in $(-\pi,\pi]$ whenever the
global input satisfies $|\im(z)|<1$.
\end{assumption}

\begin{remark}
Assumption~\ref{ass:trees} is verified numerically at 60 decimal places: zero
violations were found on a grid and 1000 random samples inside $|\im(z)|<1$.
A fully analytic interval-arithmetic proof --- propagating Im-bounds forward
through each node --- is in principle obtainable but has not been written out
in full; it remains an open verification task.
\end{remark}

By structural induction, if Assumption~\ref{ass:trees} holds for a given
compact $K\subset\{z:|\im(z)|<\pi\}$, then every composite tree built from
$\eml$, $\emlstar$, and the constant $1$ evaluates without branch-cut
crossings on $K$.

% ──────────────────────────────────────────────────────────────────────────────
\section{Consequences for the Density Theorem}
\label{sec:density}

\begin{theorem}[Stone--Weierstrass Density, Th.~4.3 of companion paper]
\label{thm:density}
Under Assumption~\ref{ass:trees}, for every compact
$K\subset\{z:\im(z)\in(-\pi,\pi)\}$, the algebra generated by
$\{\eml,\emlstar,1\}$ is dense in $C(K,\C)$.
\end{theorem}

\begin{proof}
On such $K$, Lemma~\ref{lem:branch-safety} is exact and continuous, so
$\conj{z}$ belongs to the algebra pointwise on $K$.  Combined with
Assumption~\ref{ass:trees}, the algebra contains $z$, $\conj{z}$, all
rational constants, and is closed under addition and multiplication; hence it
contains $\Q[z,\conj{z}]$.

The sub-algebra $A=\Q[z,\conj{z}]\subset C(K,\C)$ satisfies the hypotheses
of the complex Stone--Weierstrass theorem~\cite{stone1948}:
\begin{enumerate}[label=(\roman*)]
  \item \emph{Separates points}: for $z_1\neq z_2$, $z\in A$ gives
        $z(z_1)\neq z(z_2)$.
  \item \emph{Does not vanish}: $1\in A$.
  \item \emph{Closed under conjugation}: $z\in A$ and $\conj{z}\in A$ imply
        $\conj{f}\in A$ for any $f\in A$.
\end{enumerate}
Hence $A$ is dense in $C(K,\C)$.
\end{proof}

\begin{remark}
Theorem~\ref{thm:density} does not contradict Theorem~2.1 of the companion
paper (holomorphic closure of $E$): the algebra is \emph{dense} in $C(K,\C)$,
but individual $\eml$-trees remain holomorphic.  Density is an approximation
result, not a representation result.
\end{remark}

% ──────────────────────────────────────────────────────────────────────────────
\section*{Summary of Proof Status}

\begin{center}
\begin{tabular}{lll}
\toprule
Result & Status & Evidence \\
\midrule
Lemma~\ref{lem:branch-safety} (conjugation formula) &
  Proved &
  Steps 1--5 above \\
Theorem~\ref{thm:domain} (exact domain $[-\pi,\pi)$) &
  Proved &
  Section~\ref{sec:sharpness} + SymPy \\
Assumption~\ref{ass:trees} (arithmetic trees) &
  Open (analytic) &
  Numerically verified, 60~dps \\
Theorem~\ref{thm:density} (Stone--Weierstrass) &
  Conditional on Ass.~\ref{ass:trees} &
  Complete given assumption \\
\bottomrule
\end{tabular}
\end{center}

% ──────────────────────────────────────────────────────────────────────────────
\begin{thebibliography}{9}

\bibitem{odrzywolek2026}
A.~Odrzywołek,
\textit{All elementary functions from a single operator},
arXiv:2603.21852v2 (April 2026).

\bibitem{stone1948}
M.~H.~Stone,
\textit{The generalized Weierstrass approximation theorem},
Mathematics Magazine \textbf{21} (1948), 167--184.

\bibitem{terzo2008}
G.~Terzo,
\textit{Some consequences of Schanuel's conjecture in exponential rings},
Comm.\ Algebra \textbf{36} (2008), 1171--1189.

\end{thebibliography}

\bigskip
\noindent\footnotesize
\textbf{License:} MIT\quad
\textbf{Repository:} \href{https://github.com/antparis/eml_star}{\texttt{github.com/antparis/eml\_star}}\quad
\textbf{Zenodo:} \href{https://doi.org/10.5281/zenodo.19183008}{10.5281/zenodo.19183008}

\end{document}
